Democratic governance requires a population of individuals with diverse needs and
circumstances to make binding collective decisions (laws) that affect everyone.
A tax policy, zoning regulation, or healthcare mandate interacts differently with
each citizen depending on their occupation, health, location, and other personal
attributes. This interaction between shared institutions and individual
heterogeneity is a defining feature of law as a complex adaptive
system~\citep{ruhl1996fitness}, yet formal models of voting rarely account for
it.

We propose a framework that captures this structure directly. Each individual is
represented as a binary string on an NK fitness landscape~\citep{kauffman1993origins},
partitioned into a \emph{shared portion} (bits modified collectively through
voting, representing laws) and an \emph{individual portion} (bits that are
fixed and unique to each person, representing personal traits. The NK epistatic
structure couples these portions: the fitness contribution of each law-bit
depends on nearby trait-bits, and vice versa. A cross-dependency parameter
$\alpha$ controls the fraction of each bit's dependencies that cross between
portions. At $\alpha=0$, dependencies are purely internal: laws interact only
with other laws, so legislation affects all citizens identically. At $\alpha=1$,
dependencies are purely cross-portion: every law's effect depends on individual
traits, but traits lack complex internal interactions. Intermediate values produce
the richest landscapes; at $\alpha \approx 0.5$, both internal and
cross-dependencies are present, creating the highest effective complexity. This
parameter formalizes a spectrum familiar from political economy: from purely
public goods ($\alpha=0$) to policies whose effect is entirely mediated by
personal circumstances ($\alpha=1$).

Within this framework, each round of voting is a step of collective optimization
on a rugged landscape. Different voting methods: plurality, approval voting,
score voting, Borda count, instant-runoff voting (IRV), STAR voting, minimax
(Condorcet), and random dictator, aggregate individual preferences differently,
and thus navigate the landscape differently over many rounds.

This perspective suggests a useful unifying interpretation. In our model, all
voting methods receive the same input: a utility matrix
$U \in \mathbb{R}^{M \times 2^V}$ whose $(i,p)$ entry is the fitness voter~$i$
would obtain under proposal~$p$. A voting rule can therefore be viewed as a
\emph{transformation of the same utility field}: different methods preserve,
compress, or discard different aspects of $U$ before selecting a winner. For
example, plurality keeps only each voter's top choice, score voting sums raw
cardinal values, and Borda discards magnitudes but preserves full within-voter
ordinal structure. From this viewpoint, voting methods differ not only in
normative properties but also in the kind of information processing they
perform.

Beyond direct democracy, we develop a unified model of \emph{representative
democracy} grounded in two longstanding debates in political theory. The first is
the \emph{trustee--delegate debate}: does a representative exercise independent
judgment (Burke's trustee~\citep{burke1774bristol}) or faithfully implement
constituent preferences (the delegate model)? We parameterize this as
$p_{\mathrm{self}} \in [0,1]$, the probability a candidate's platform reflects
their own interest rather than their constituency's collective welfare. The second
is the question of \emph{identity versus policy voting}: do citizens elect
representatives who share their background and circumstances (descriptive
representation~\citep{mansbridge1999should}) or those whose platforms promise the
best outcomes? We parameterize this as $\beta \in [0,1]$, the weight on
identity-based versus policy-based candidate evaluation. Candidates' constituencies
are defined by nearest-neighbor Hamming distance in individual traits. Those who
share your circumstances are your representatives. The cross-dependency $\alpha$
then plays a natural role: when $\alpha=0$, individual traits are irrelevant to
policy outcomes and identity voting carries no information; when $\alpha=1$,
traits fully mediate policy effects, making identity a strong proxy for policy
alignment.

We study two core questions. First, how does voting method choice affect long-run
collective welfare (mean fitness, variance, and distributional outcomes) across
varying landscape ruggedness $K$ and law-trait coupling $\alpha$? Second, how
does the introduction of representation, and the choice of $\beta$ and
$p_{\mathrm{self}}$, alter these dynamics, and in what parameter regimes does
representation help or hurt?

Our main finding is that voting methods undergo sharp phase transitions as a
function of landscape complexity: total score voting dominates on smooth
landscapes, a generalized scoring rule with $p=0.35$ at low-to-moderate
ruggedness, Borda count across a wide middle range, and STAR voting at the
highest complexity. These transitions are
well-described by a two-parameter empirical formula that reduces the $(K, \alpha)$
plane to a single complexity axis for visualization. Borda count is the most consistent method,
achieving both the highest mean fitness and lowest variance across the majority
of the parameter space. Under representative democracy, the complexity-dependent
regime structure is reshaped even under favorable conditions: total score voting
dominates across most regimes, and as identity voting increases, plurality
performs surprisingly well while sophisticated methods lose their advantage.
Conceptually, these results also suggest a
new research question: whether voting rules can be ordered by computable
properties of the utility transformation they implement, and whether such
properties predict performance across environments.
